\section{Observable and source model}\label{sec:cf-framework}
A frequently used schematic representation is the Koonin--Pratt integral
\begin{equation}\label{eq:cf-kp}
C(\bm{k}^{*})=\int \dd^3r^{*}\,S(\bm{r}^{*})\,
\abs{\Psi_{\bm{k}^{*}}^{(-)}(\bm{r}^{*})}^2,
\end{equation}
with a normalized pair-separation distribution $\int\dd^3r^*\,S(\bm r^*)=1$. This is a schematic stationary-pair form; spin averaging, normalization, channel coupling, and experimental kinematic selection must be specified before quantitative use. Historical starting points include \textcite{Koonin1977,Pratt1986} and an analytical treatment of final-state effects by \textcite{Lednicky1982}.

\subsection{Frames and symmetrization}
An analysis must say whether the coordinates are in the pair rest frame or a longitudinally comoving frame, distinguish $k^*$ from momentum difference $q$, and include symmetrization or antisymmetrization when identical particles are measured. Spin populations and detector acceptance can change the apparent low-$k^*$ structure.

\subsection{Physics versus reconstruction}
An observed pair distribution is not simply $C(k^*)$: tracking efficiency, close-track rejection, merging/splitting, particle identification, momentum resolution and event-mixing normalization can all reshape it. Comparisons should forward-fold a theoretical prediction through these effects or unfold with a validated response model.
